\section{Experimental Setup}
\label{sec:setup}

\subsection{Benchmarks}

We assemble a six-benchmark suite that spans knowledge, commonsense, mathematical reasoning, and logical reasoning, with a mix of option counts and answer formats. TruthfulQA-MC1 \citep{lin2022truthfulqa} contributes adversarial factual questions with four or five options; we shuffle option order per question to avoid position bias. PIQA \citep{bisk2020piqa} contributes two-choice physical commonsense problems. LogiQA \citep{liu2020logiqa} contributes four-choice logical reading comprehension. LSAT logical reasoning and SAT mathematics, both from the AGIEval suite \citep{zhong2023agieval}, contribute five-choice analytical reasoning and four-choice quantitative items respectively. GSM8K \citep{cobbe2021training} contributes free-form grade-school mathematics problems answered with chain-of-thought reasoning and a final numeric value. We sample fifteen questions per benchmark uniformly at random with a fixed seed, giving ninety questions in total. The mixture is deliberately diverse in difficulty: as Section~\ref{sec:results} shows, the strongest models in our pool saturate five of the six benchmarks while all models find LogiQA and, for some models, GSM8K and SAT mathematics challenging, which supplies both the degenerate and the high-entropy regimes that our analysis requires.

\subsection{Models}

We study six commercial models from four families, accessed through three providers under their free tiers: Qwen3.8-Max and Qwen3.8-Omni-Flash from the Qwen family \citep{qwen2024qwen}; Agnes3.0-Flash and Agnes2.5-Flash from the Agnes family; GPT-6-Luna from the GPT family \citep{openai2023gpt4}; and Gemini3.8-Flash from the Gemini family. The models span a range of capability tiers and post-training styles, including one model that emits an explicit private reasoning trace and one multimodal variant. All models are queried through OpenAI-compatible endpoints between October 16 and October 19, 2026. Sampling uses temperature $0.7$ with default nucleus probability; greedy calls use temperature $0$. For the two Qwen models we disable the proprietary thinking mode on multiple-choice tasks to control token expenditure and enable it on GSM8K, where chain-of-thought is necessary; all other models run in their default configurations, and every model receives the identical instruction format described below.

\subsection{Prompt design and paraphrase generation}

The canonical prompt for a multiple-choice question consists of the passage when present, the question stem, the option list labeled $(A)$ through $(E)$, and a fixed instruction that requires a direct answer in the form \texttt{Answer: <letter>} without explanation. GSM8K prompts instead require step-by-step reasoning ending with \texttt{Answer: <number>}. Direct answering on multiple-choice tasks keeps the comparison uniform across models, keeps token costs bounded, and isolates the aggregation question from chain-of-thought length effects; the GSM8K condition covers the reasoning-trace regime.

Paraphrases target the parts of the prompt where semantic-preserving variation is well defined: the question stem for TruthfulQA, LSAT, and LogiQA, and the instruction wording for GSM8K, PIQA, and SAT mathematics, whose stems contain problem-specific content such as numbers that should not be altered. This design yields two paraphrase loci, and we treat the locus as a factor in the analysis. Stem paraphrases are generated by a seventh model (Qwen3.7-Flash) that is not part of the evaluation pool, prompted to rewrite each stem in six ways while preserving all facts, entities, numbers, and the question type. Instruction paraphrases are six handcrafted variants per dataset that vary wording while keeping the literal answer-format requirement fixed, which keeps answer extraction reliable.

Every generated paraphrase passes a three-stage validation filter. A lexical stage rejects paraphrases that are identical to the original or whose length ratio falls outside $[0.4, 2.5]$. A semantic stage embeds the original stem and each candidate with a text embedding model and rejects candidates whose cosine similarity to the original falls below $0.75$. A form stage, applied to factual questions, rejects rewrites that change the question type, for example from a wh-question to a yes-no question. Candidates that fail any stage are replaced by regenerated candidates until six valid paraphrases per question are obtained; the final set achieves a mean embedding similarity of $0.86$ with the originals.

\subsection{Execution and quality control}

The experiment executes the sixteen-call matrix of Section~\ref{sec:method} for every question and model through an asynchronous runner with per-provider rate limiting, automatic retry with exponential backoff, and append-only checkpointing. Answer extraction uses the mandated \texttt{Answer:} format first and falls back to scanning the final output lines for a well-formed option letter or number, with the option range constrained by the question. Every stored record includes the extracted answer, its correctness, token usage, and a content snippet for auditing. Extraction succeeds on over 96\% of calls for every model reported in this paper; failed extractions are excluded from the corresponding question's conditions, and questions with any missing cell are excluded from that model's analysis, so all reported accuracies are computed over complete question matrices. Per-model completed question counts, which vary between 25 and 90 owing to provider rate limits, are stated alongside every result, and pooled analyses use the intersection of completed questions across the models involved.

\subsection{Conditions and metrics}

Each model and question yields the greedy baseline, $\mathrm{SC}\text{-}k$ and $\mathrm{PE}\text{-}k$ for $k \in \{2, 4, 6\}$, the hybrid $\mathrm{HB}\text{-}6$, the combined $\mathrm{COMB}\text{-}12$, the agreement statistics $c_{\mathrm{S}}$ and $c_{\mathrm{P}}$, the normalized answer entropy, and the fragility indicator. Accuracy is the fraction of questions answered correctly. Error detection is evaluated by AUROC, AURC, and E-AURC as defined in Section~\ref{sec:method}. Budget curves plot accuracy against $k$ for both pure strategies. All analysis code operates on the raw prediction logs, and we release the logs together with the prompts and paraphrases.
